% GENERATED by md2tex.py from sections/07_discussion_limitations.md - edit the .md, not this file
\section{Discussion and Limitations}

\textbf{Scale scope of the capability legs.} The rank law and the recall
separation are small-model, synthetic-task results by design: the
encoder family trains in minutes and the head-to-head contender is a
14M-parameter model on a controlled binding task. The strength of that
regime is causal closure (train-time rank forcing, state zeroing,
frozen margins); its limit is external validity. Nothing here shows the
recall separation surviving at language-model scale or on natural
data, and the matched-budget caveat bounds every transformer
comparison: a differently scaled or differently trained transformer
requires its own matched re-run before any extrapolation. The
learning-rate search for the transformer arm on the recall task itself
has since run. Under the identical frozen protocol, a grid of four
learning rates spanning $10^{-4}$ to $3 \times 10^{-3}$ (three seeds
each, 20,000 matched steps) reads at or below chance at every point:
mean $\mathrm{acc}_A$ is 0.0142 at $\mathrm{lr}=10^{-4}$, 0.0283 at
$3\times10^{-4}$ (the shared default used elsewhere in this paper),
0.0320 at $10^{-3}$ (the best of the grid, $1.02\times$ chance), and
0.0283 at $3\times10^{-3}$; 0 of 12 individual (learning rate, seed)
cells and 0 of 4 grid means clear the frozen 0.09375 demonstration bar
($3\times$ chance at $K=32$). Per the
pre-registered decision rule, the transformer now joins the
parameter-matched vector-state ablation as a second baseline that
fails to demonstrate recall after an explicit learning-rate search;
the ablation comparison remains the registered verdict carrier
(Section 4.2), and this measurement changes no registered tier.

The grid also closes the leading unexcluded explanation named in
Section 6 for the chance-level reading: the best-optimizing learning
rate ($10^{-4}$; final training loss 6.55 to 6.60, well below the
shared default's approximately 7.5) reads $\mathrm{acc}_A = 0.0142$,
6.3 standard deviations below chance, while the two worse-optimizing
learning rates ($10^{-3}$ and $3\times10^{-3}$; final loss
approximately 7.7 to 7.8) read closer to chance. Optimization quality and recall dissociate on
this arm: giving the transformer the learning rate that best optimizes
the language-modeling objective does not surface recall, so the
chance-level reading is not attributable to an under-optimized
transformer. This closes the specific unexcluded explanation named
above; it does not establish an architectural impossibility. The
matched-budget caveat still bounds the reading: this is a two-layer,
approximately 15M-parameter, 20,000-step, from-scratch regime, and the
MQAR and induction-head literature's positive transformer results
(Section 6) live at other scales, architectures, and budgets.

\textbf{One stress point is a locate-only null.} At the disclosed
above-capacity stress load ($K/d = 0.75$), the completed three-arm
table reads chance in every arm: contender 0.0189, ablation 0.0195,
transformer 0.0218, against chance 0.0208. The
table is locate-only by pre-registration, not a verdict-grade cell; no
capability claim is made at this load, and it bounds the single-hop
recall separation to loads below this stress point at this training
budget.

\textbf{Scope of the pathology leg.} The ladder result is a geometry claim:
span fraction worsens monotonically with scale, but this paper does not
tie the drift to a downstream performance cost, and validation loss is
neutral in every mitigation cell measured. The tie between the geometric
drift and any behavioral cost is the outstanding scientific question,
and the 392M token-budget confound (Section 5.3) caps what the
transfer wave can say across scales.

\textbf{Multi-hop composition is open.} The two-hop task is
verdict-INDETERMINATE, with partial recall in three of nine pooled
contender seeds, none generalizing to held-out hop depths (Section
4.5), and a separate compositional-depth program in this research line
is at the design-and-calibration stage with a pre-registered lens
question resolved but no model verdict; we claim neither. Depth
generalization is future work, stated as such.

\textbf{Instrument validity as a recurring risk.} Three times in this
program (Sections 2.5, 3.3, 4.3), the first instrument pointed at a
true phenomenon read false. In each case the fix was minor once the
defect was seen. Representational claims about matrix
states need independent checks of the kind Section 2.5 catalogs,
because errors in either direction arise from instrument choice alone
at this scale: the wrong-tap probe read a real capability as absent,
and the identity-padded target sold ambient structure as recovery.

